\section{Some examples of Dynamical Systems results}\label{sec:some-exampl-dynamical-sys-res}

A very first example of a result in Dynamical Systems might be the following classical

\begin{theorem}[Banach fixed point theorem]\label{thm:Banach-fixed-point-thm}
  Let $(M,d)$ be a complete metric space and $f\colon M\carr$ be \textbf{uniform contraction,} \ie~$f$ is a Lipschitz function and its \textbf{Lipschitz constant}
  \begin{displaymath}
    \lambda=\Lip(f):=\sup_{x\neq y}\frac{d\big(f(x),f(y)\big)}{d(x,y)}< 1.
  \end{displaymath}

  Then there exists a unique point $p\in M$ such that $f(p)=p$ and it holds
  \begin{displaymath}
    \lim_{n\to+\infty} f^{n}(x)=p, \quad\forall x\in M.
  \end{displaymath}
\end{theorem}

\begin{proof}
  Given any $x\in M$, one can show by induction on $n$ that
  \begin{displaymath}
    d\big(f^{n}(x),f^{n+1}(x)\big)\leq \lambda^{n}d(x,f(x)), \quad\forall n\in\N_{0}.
  \end{displaymath}
  So, for every $m,n\in\N$ it holds
  \begin{displaymath}
    \begin{split}
      d\big(f^{m}(x),f^{m+n}(x)\big)&\leq \sum_{i=0}^{n-1} d\big(f^{m+i}(x),f^{m+i+1}(x)\big) \\
      &\leq \sum_{i=0}^{n-1}\lambda^{m+i}d\big(x,f(x)\big) \leq \frac{\lambda^{m}}{1-\lambda} d\big(x,f(x)\big).
    \end{split}
  \end{displaymath}
  Hence, the sequence ${\big(f^{n}(x)\big)}_{n\geq 0}$ is Cauchy, and since $M$ is complete, it converges to a point $p\in M$. By continuity of $f$, $p$ is a fixed point of $f$. On the other hand, since $f$ is a contraction, it has at most one fixed point.
\end{proof}

\begin{exercise}[Perturbations of contractions]
  Let $f\colon\R^{d}\carr$ be a uniform contraction and $r>0$ be arbitrary. We define the following balls
  \begin{displaymath}
    B_{r}^{0}(f):=\left\{g\in C^{0}(\R^{d},\R^{d}) : \norm{f-g}_{C^{0}}:=\sup_{x\in\R^{d}}\norm{f(x)-g(x)}<r\right\},
  \end{displaymath}
  and
  \begin{displaymath}
    B_{r}^{L}(f):=\left\{g\in C^{0}(\R^{d},\R^{d}) : \norm{f-g}_{L}:=\norm{f-g}_{C^{0}} + \Lip(f-g)<r\right\}.
  \end{displaymath}

  Show that given any $r>0$, there is a $g\in B_{r}^{0}(f)$ which is not a contraction and such that $\Fix(g):=\{x\in\R^{d} : g(x)=x\}$ is an infinite set.

  On the other hand, show that given the uniform contraction $f$, there exists $r_{0}>0$ such that every $g\in B_{r_{0}}^{L}(f)$ is a uniform contraction, and hence it exhibits a unique fixed point $p_{g}\in\R^{d}$. Finally, prove that given any $\eps>0$, there exists $\delta\in(0,r_{0})$ such that
  \begin{displaymath}
    \norm{p_{f}-p_{g}}<\eps, \quad\forall g\in B_{\delta}^{L}(f).
  \end{displaymath}
\end{exercise}

Another important example is the following one:

\begin{example}[Gradient flows]\label{exam:gradient-flows}
  Let $M$ be the unit sphere $\bb{S}^{2}:=\{v\in\R^{3} : \norm{v}=1\}$ and consider the function $L\colon\bb{S}^{2}\to\R$ given by
  \begin{displaymath}
    L(x,y,z):=-z, \quad\forall (x,y,z)\in\bb{S}^{2},
  \end{displaymath}
  and the flow $\Phi\colon\R\times\bb{S}^{2}\to\bb{S}^{2}$ generated by the vector field $X(v):=-\nabla L(v)$, for every $v\in\bb{S}^{2}$ and where $\nabla$ denotes the gradient operator on $\bb{S}^{2}$.

  One can easily show that $(0,0,-1)$ and $(0,0,1)$ are the only singularities of $X$ (\ie~the critical points of $L$), or equivalently, the fixed points of $\Phi$, and it holds
  \begin{displaymath}
    \Phi^{t}(v)\to (0,0,-1), \quad\text{and}\quad \Phi^{-t}(v)\to (0,0,1),
  \end{displaymath}
  as $t\to+\infty$, for every $v\in\bb{S}^{2}\setminus\{(0,0,\pm 1)\}$.

  Generalizing this example, one can consider an arbitrary compact differentiable manifold $M$ endowed with a Riemannian structure $\langle\cdot,\cdot\rangle$ and a smooth function $L\colon M\to\R$. Then we can define the gradient of $L$ as the only vector field $\nabla L$ such that
  \begin{displaymath}
    YL = dL(Y) =\langle\nabla L, Y\rangle,
  \end{displaymath}
  for every vector field $Y$ on $M$.

  So, if we consider the flow $\Phi\colon\R\times M\to M$ induced by the vector field $X:=-\nabla L$, one can easily check that
  \begin{displaymath}
    \partial_{t} (L\circ\Phi^{t})(p)\big|_{t=0}=-\norm{\nabla L(p)}^{2} = -\langle\nabla L(p),\nabla L(p)\rangle,\quad\forall p\in M.
  \end{displaymath}

  That means, given any point $p\in M$ the map $\R\ni t\mapsto L\circ\Phi^{t}(p)\in\R$ is not increasing, and it is strictly decreasing if and only if $p$ is a regular point for $L$ is a regular point of $L$, \ie~$dL_{p}\colon T_{p}M\to\R$ is not identically zero.

  Moreover, if $L$ has a finite number of critical points (\ie~points where the derivative $dL$ vanishes), then one can show that for every $p\in M$, there exist critical points of $L$ $\alpha_{p},\omega_{p}\in M$ such that
  \begin{displaymath}
    \Phi^{-t}(p)\to\alpha_{p}, \quad\text\quad \Phi^{t}(p)\to\omega_{p},
  \end{displaymath}
  as $t\to+\infty$.
\end{example}
